\section{A direct consistency quartic in 576 unrestricted integers}
\label{sec:direct576}

The direct proof-certificate route admits a new simplification: a
quadratic bound on a proposed power output can replace its sign
predicate.  Once the indexed Pell coordinates are oriented, a small
signed residue identifies the output and recovers its sign.  Combining
this with two local source identities and a factorial digit obstruction
gives the following explicit quartics.

\begin{center}
\begin{tabular}{lrrrr}
\hline
Equation & Variables & Quadratic rows & Terms & Height\\
\hline
Previous principal &604&372&9,676&2,445,790,974,704\\
New principal &\textbf{576}&368&\textbf{7,172}&2,445,790,974,704\\
Previous lower height &609&377&8,823&964,213\\
New lower height &\textbf{582}&374&\textbf{6,200}&964,213\\
\hline
\end{tabular}
\end{center}

All four total degrees are exactly four.  The two new equations are
literally
\[
 P_{576}(x)=\sum_{j=0}^{367}R_j(x)^2=0,
 \qquad
 P_{582}(x)=\sum_{j=0}^{373}\widetilde R_j(x)^2=0,
\]
where the accompanying JSON files list every quadratic residual and
every nonzero expanded coefficient.  There are no coefficient ports or
unmaterialized variable powers.  The first equation saves 28 variables
and 2,504 monomials at unchanged height; the second saves 27 variables
and 2,623 monomials.  These are measured improvements, not a minimality
claim.  This direct C222 route does not use K215, Jones compression or
Sun's integer reduction.

The precise new theorem is integer-zero existence equivalence with the
pinned 75-predicate direct-C222 source.  Its existing ordinary
refutation/consistency calibration transfers separately.  All arguments
here are finite integer arithmetic that can be performed in ZFC; no
consistency or external soundness assumption is made.  A new Lean,
Rocq or sealed object-language ZFC proof object is not claimed.

\subsection{The exact source change}

The source is pinned at ProveIt commit
\texttt{84e16e43092614bae262ca59651cb9d77e31a255}, in the
\texttt{factorial-radius-output-sign-profile/raw} research packet.
Its body has 1,259 coordinates and 1,177 quadratic residuals.  A
predicate $E\ge0$ is represented by
\[
 E-a^2-b^2-c(c+1)=0
\]
with three private integer variables.  These witnesses exist exactly
when $E\ge0$, by the three-square theorem applied to $4E+1$.

Delete exactly the twelve scalar predicates with raw indices
\[
 6,104,106,575,584,592,600,609,618,627,636,825.
\]
Keep the other 63 predicates.  Put
\[
 S=\{6,106,573,583,591,599,607,617,625,635\}.
\]
Change only body residual 1176, originally $x_{643}-x_{1258}$, to
\[
 x_{643}-x_{1258}+\sum_{i\in S}(x_i-x_i^2).
\]
This adds no coordinate and preserves quadratic degree.  Every other
body residual is unchanged.  The resulting raw equation has 1,448
integer coordinates and 1,240 quadratic residuals.

Let $\mathcal A$ be the set of all 43 distinct base, exponent and
output indices in the 21 actual power calls.  The original frame
gate identities now give exactly
\begin{equation}
 W=2+\sum_{i\in\mathcal A\setminus S}x_i+\sum_{i\in S}x_i^2,
 \qquad W=x_{643}.
 \label{eq:direct-signed-frame}
\end{equation}
The complete list $\mathcal A$ and the actual source-row ideal are
included in the certificate.  A squared contribution bounds either
sign of an integer because $z^2\ge |z|$.

\subsection{Signs and a dominating Pell frame}

Write $f=x_{87}$, $r=x_{94}$, $s=x_{95}$ and $N=x_{93}$.  The
retained predicates give $f\ge2$, $r,s\ge0$.  Exact source rows
78--85 imply
\[
 N=fr+(f-1)s+f\ge2.
\]
Separately, an eleven-row polynomial identity gives
\begin{equation}
 f(x_{104}+1)=f x_{59}+x_{131}+3.
 \label{eq:direct-output104}
\end{equation}
The two right-hand coordinates retain direct nonnegativity predicates.
Thus $x_{104}+1>0$ and $x_{104}\ge0$, without any power interpretation.
The full eleven integral polynomial multipliers are supplied.

Filtering the old prefix proof now proves 604 of the 643 semantic
coordinate signs.  Its only postponed special bound concerns $x_{130}$.
The only unknown-sign frame arguments are precisely the ten members
of $S$.  All actual bases and seed exponents are already positive.
In particular, with $t=x_4$, $K=x_{86}$ and $f_1=x_{107}$,
\begin{align*}
 &t\ge1,\quad K\ge3,\quad f_1\ge0,\\
 &M=x_{129}=1+f(x_{59}+1)\ge3,\\
 &b=x_{530}=M(1+N+x_{113})\ge9,\qquad x_{531}=b+1\ge10,\\
 &x_{103}=1+(t+1)f\ge5.
\end{align*}
Consequently the new frame gives, for every actual call $(n,k,m)$,
\begin{equation}
 W\ge3,\qquad 1\le n,k<W,\qquad |m|<W.
 \label{eq:direct-frame-bounds}
\end{equation}

The retained initial modulus and its identity
$T_0=2A(t+1)-(t+1)^2-1$ force $A=x_{644}\ge2$.
The unchanged bootstrap norm is
\[
 A^2-1=W^3(W+2)Z^2.
\]
Here $Z\ne0$.  At Pell base $W+1$ this is the solution
$(A,W|Z|)$.  Classification and the recurrence modulo $W$ imply
that its positive index is a multiple of $W$.  Hence
\begin{equation}
 A>W^W,
 \qquad T_n=2An-n^2-1\ge A>W^W.
 \label{eq:direct-large-modulus}
\end{equation}
The last inequality follows from the exact decomposition
\[
 T_n-A=(2n-1)(A-n-1)+(n-1)(n+2).
\]

The retained $V\ge1$ predicate and the unchanged product identity
$V=(\prod_iY_i^2)(v+1)$ make every seed ordinate nonzero and give
$Y_i^2\mid V$.  The common norm implies $|U|\ge A$.  From
$B=1+4V(q+1)$ and $B\equiv A\pmod U$, the case $q=-1$ is impossible:
it would force $U\mid A-1$, although $0<A-1<|U|$.  Every other $q$
gives $|B|\ge3$.  Thus the pinned signed-base index theorem is
available whenever the particular short-interval condition
$k\le |Y|$ has been established.

\subsection{A power output may initially be signed}

\paragraph{Signed-output lemma.}
In this section $\chi_k(a),\psi_k(a)$ use an index subscript:
they start at $(1,0)$ and satisfy
$\chi_{k+1}=a\chi_k+(a^2-1)\psi_k$,
$\psi_{k+1}=\chi_k+a\psi_k$.
Assume \eqref{eq:direct-frame-bounds}--\eqref{eq:direct-large-modulus},
$n\ge2$, and
\[
 |X|=\chi_k(A),\quad Y=\psi_k(A)>0,\qquad
 X-(A-n)Y-m=T_nh,
\]
where $m,h$ are unrestricted integers.  Then
$X=\chi_k(A)$ and $m=n^k$.

\paragraph{Proof.}
Put $p=\chi_k-(A-n)\psi_k$ and
$q=\chi_k+(A-n)\psi_k$.  The Pell recurrence and norm give
$p\equiv n^k\pmod{T_n}$ and $pq=1+T_n\psi_k^2$.
If $X<0$, then $T_n\mid mn^k+1$.  This dividend is nonzero
for every integer $m$, since $n^k\ge2$, but
\[
 |mn^k+1|\le(W-1)^W+1<W^W<T_n.
\]
This is impossible.  Hence $X>0$, and $T_n\mid m-n^k$.  Now
\[
 |m-n^k|\le W-1+(W-1)^{W-1}<W^W<T_n,
\]
so $m=n^k$.  Notice that positivity of $Y$ is essential:
if both main coordinates could be negative, $m=-n^k$ would survive.
Every application below first supplies that positive ordinate.

\subsection{Restoring the deleted predicates in order}

\paragraph{The eight range quotients.}
For actual calls $10,12,13,14,15,17,18,20$, all eight original
$Y\ge k$ predicates remain.  The respective outputs are
$573,583,591,599,607,617,625,635$; all bases are $x_{531}\ge10$.
The signed index theorem followed by the signed-output lemma makes
these outputs their intended positive powers.  Their exact range
rows are
\[
 m=(Me)q+r,\qquad e=r+s+1,\qquad r,s\ge0.
\]
Here $M=x_{129}$, not the radix $b=x_{530}$.  Since
$0\le r<e\le Me$, a negative $q$ would make $m<0$.
All eight quotients and their 24 dependent signs are recovered.
Only the seven prefix signs
\[
 6,106,124,125,130,132,133
\]
remain unproved.

\paragraph{The first factorial.}
None of those seven coordinates is a sign premise of the pinned
local first-factorial argument.  Put $d=fK^2$ and $C=x_{91}$.
Its exact inputs are
\[
 C=r+s+1,\quad fC=N+s,\qquad
 x_{532}=x_{533}(C+ba)+u,
\]
\[
 x_{533}=u+v+1,\qquad b=C+z+1,
\]
with $a,u,v,z\ge0$, together with actual calls $1,5,6$.
All three outputs $N,x_{532},x_{533}$ remain known nonnegative.

For either missing first-factorial index, $|Y|\ge2$ implies
$|Y|\ge2A>W>k$, so the index theorem and a nonnegative-output
supporting call recover its true positive pair.  If $|Y|=1$, the
four possible residues $n,-2A+n,2A-n,-n$ and the small-output bound
leave exactly $X=A,Y=1,m=n$.  Thus the effective indices are
$e\in\{K,1\}$ and $j\in\{K+d,1\}$, with
\[
 N=(d+1)^e,\quad x_{533}=b^e,\quad x_{532}=(b+1)^j.
\]
If $j=1,e=K$, the digit equation contradicts $C\ge1$; if
$j=e=1$, it forces $C=1$ and then $f\ge d+1>f$.
Therefore $j=K+d$.  If $e=1$, the bound
$b\ge(f+1)(d+2)>K+d$ identifies the first binomial digit as
$C=K+d$, and the ceiling equation forces $f=1$, again impossible.
Both index guards are restored.  The pinned local factorial lemma
now gives $f=K!$, so $f$ is even.  Its complete low-digit and
self-referential-radius proof is bundled; it uses these three calls,
not an interpretation of the remaining powers.

\paragraph{Orienting the t-pair with a different output.}
The base $x_{103}=1+(t+1)f$ is odd.  The t-seed norm and its
actual call-2 residue modulo two imply $x_{104}$ is odd, without
interpreting t.  Hence $t$ and $k_{122}=t+x_{104}+2$ have opposite
parity.  An even index excludes $|Y|=1$: the signed abscissa
half-period argument makes the comparison index odd, while its
ordinate congruence modulo four makes it even.  Thus its true index
is recovered, and a positive-output call orients its pair.

If t is even, use \emph{call 2}, whose output $x_{104}\ge0$ was
proved in \eqref{eq:direct-output104}; its base is at least five.
If $k_{122}$ is even, use \emph{call 9}, with output $x_{560}\ge0$.
The retained product of these two index slacks then makes the
other slack nonnegative.  The indexed t-pair is oriented by call 2
also when t is odd.  This change of supporting call avoids assuming
the missing sign of call 0's output $x_6$.

With that positive t-pair, the signed-output lemma at bases $t+1$
and $f$ recovers $x_6=(t+1)^t>0$ and $x_{106}=f^t>0$.
The remaining source rows give
\[
 x_{125}=x_{123}x_{107}x_{106}\ge0,\quad
 x_{125}=M+x_{126}x_{130},\quad
 x_{126}=M+x_{131}+1>M.
\]
Consequently $x_{130}\ge0$; its dependent signs follow.  The
entire 643-coordinate semantic prefix is now natural.

\paragraph{The second factorial's shifted index.}
It remains to recover the deleted predicate at 825.  With t exact,
actual calls 4 and 8 give
\[
 N_1=x_{113}=(1+f_1t^2)^t\ge1,\qquad x_{547}=b^t.
\]
The second factorial rows give
$C_1=x_{114}+x_{115}+1\ge1$ and $f_1C_1=N_1+x_{115}$,
so $f_1\ge1$.  Its shifted seed has index $t+f_1t^2$.
The large-ordinate case recovers its true positive pair as above.
In the small case $|Y|=1$, the signed residue analysis forces
$x_{546}=b+1$.  Its exact digit equation is
\[
 x_{546}=b^t(C_1+ba_1)+u_1,\qquad
 C_1\ge1,\quad a_1\ge0,\quad0\le u_1<b^t.
\]
For $t\ge2$ the right side is at least $b^t>b+1$.
For $t=1$, division forces $C_1=1$, whereas
$f_1C_1=f_1+1+x_{115}>f_1$.  Both cases are impossible.
Thus the omitted shifted-index slack is strictly positive.

\subsection{Completeness requires a fresh private frame}

All 75 parent predicate values have now been recovered.  The unchanged
Pell and small-residue arguments give every true power value; they
use only the dominating bounds, not the literal old linear frame.
Nevertheless, adjoining guard triples alone does \emph{not} produce
an old source zero: its original frame equality usually fails.

Hold the natural semantic prefix fixed and reconstruct the private
Pell data.  Choose either its old linear frame $W_0$ or the new squared
one, and set
\[
 A_0=\chi_{W_0}(W_0+1),\qquad
 Z_0=\psi_{W_0}(W_0+1)/W_0.
\]
The recurrence modulo $W_0$ proves the division is exact.  Choose
all main pairs at their actual indices.  In the finite quadratic
ring modulo $L=4\prod_iY_i^2$, the unit
$A_0+\sqrt{A_0^2-1}$ has finite order.  This supplies a positive
common Pell solution with $U\equiv1$, $V\equiv0\pmod L$.
The norm gives $\gcd(U,V)=1$, and $U$ is odd.  CRT supplies
$B\equiv A_0\pmod U$, $B\equiv1\pmod{4V}$, with $B>A_0$.
Take the comparison pairs at their true indices and use their
integral progression quotients.  The product quotient is
$V/(\prod_iY_i^2)-1\ge0$.

For each power call, its integral quotient follows from
\[
 H_0=1,\quad H_{j+1}=nH_j+T_n\psi_j(A_0),\quad
 H_j=\chi_j(A_0)-(A_0-n)\psi_j(A_0).
\]
Reevaluate the derived t+1 pair, discriminant, moduli and private
arithmetic gates.  Every one of the 109 private nongate residuals
is the sole nongate, with coefficient $\pm1$, in an exact row
identity for a target just satisfied.  All other terms are unit
gate equations.  The independent audit checks this entire coverage,
including the 474 private topological gate definitions.  Fresh
three-square witnesses satisfy every private guard.

Using the old linear frame constructs an old 75-predicate zero from
any new zero.  Conversely, use the old source's true powers and the
same completion with the squared frame, omitting the twelve predicates.
This constructs a new zero.  The maps establish existence equivalence;
they are not asserted to be inverse on private witnesses.

\subsection{Exact algebraic verification}

The final projections contain respectively 872 and 866 monic
eliminations.  Each deleted coordinate has a unique polynomial
integer restoration, every surviving residual stays quadratic,
and no unconstrained variable is discarded.  Reversing the sequence
gives a bijection of complete integer zero sets, including guard rows.
An independently authored checker verifies the source changes,
604-sign initial dependency boundary, all staged interfaces and
all 109 private completion targets.  A separate checker imports no
optimizer, replays every signed restoration and collects every
quartic coefficient.  Both normal and optimized Python modes pass,
and changed source rows, guards, inverses and coefficients are rejected.

These finite checks establish the literal polynomials and their
algebraic transformations.  The unbounded arguments are the Pell,
integer-order, local factorial and finite CRT proofs given above and
in the bundled pinned sources.  They do not prove Con(ZFC) itself.
The next concrete research questions are to minimize the outputs
needed to orient all seed pairs, and to optimize the signed frame
together with degree-preserving monic projection: a square is free
in the raw coordinate count but may prevent a later elimination.
